\section{Mesopotamian}

This is the oldest written cosmology we have. It comes from the land between the two rivers, the Tigris and the Euphrates. It underlies much of what follows. Long before the later scriptures of the region, the scribes of Sumer, Akkad, Babylon, and Assyria already told of a watery beginning, a council of gods, and a humanity made of clay to carry the work. This is the \textbf{pre-Judaic layer}, the root from which many later images grew.

\subsection{The Creation Model}

In the beginning there was only water. Two waters mingled together. One was \textbf{fresh}, calm and sweet. The other was \textbf{salt}, vast and churning. Where the two waters met and ran into each other, the first gods were born.

The young gods were loud. Their noise broke the rest of the deep. The fresh water resolved to kill them and have peace again. But the cleverest of the young gods learned the plot first. He put the fresh water to sleep and slew it, and he built his own house upon the slain body.

Then the salt water rose in grief and rage. She bred an army of monsters and set a champion over them, and she gave that champion the warrant of supreme rule. The gods were afraid. None would face her. At last a young storm-god stepped forward. He asked for kingship over all the gods as his price, and they granted it.

He went out against the chaos-mother with the winds and a net and a bow. He drove the winds into her open mouth so she could not close it, and he loosed his arrow and split her. Then he cut her body in two. From the upper half he made the \textbf{sky}, and from the lower half the \textbf{earth}. From her eyes he set the great rivers flowing. He fixed the stars and the moon and the sun in their courses and laid down the calendar.

The gods were still weary, for someone had to do the digging and the work. So the champion took the blood of the rebel champion of chaos, mixed it with \textbf{clay}, and made \textbf{humanity}. Humans were made to carry the labor of the gods and to feed them with offerings. Then the gods built a great city for their king and proclaimed his many names.

\subsection{The Model with Its Names}

Two cosmogonies sit side by side in the tradition. The older \textbf{Sumerian} account begins from a sea-mother. The later \textbf{Babylonian} account, the \textbf{Enuma Elish}, begins from a primordial pair. Both agree that water is first and that humans are made to serve.

\subsubsection{The Sumerian Layer}

The Sumerian picture is not one continuous epic. It is reconstructed from many texts and god-lists. The sequence is birth, then separation, then ordering.

\begin{itemize}
 \item \textbf{Nammu} is the primeval sea, the watery abyss before all else, the mother who gives birth to heaven and earth.
 \item From Nammu come heaven (\textbf{An}) and earth (\textbf{Ki}), joined as one undivided mass called \textbf{An-Ki}.
 \item Their union produces \textbf{Enlil}, the god of air and wind.
 \item \textbf{Enlil separates heaven from earth.} An carries off the sky, Enlil carries off the earth, and the air between them becomes the room of the world.
 \item \textbf{Enki}, the god of wisdom and the sweet underground waters, then orders the world. He fills the rivers, sets the marshes and fields, and assigns each god its station.
\end{itemize}

Humanity is made because the lesser gods grow tired of digging the canals. Nammu wakes Enki, and the mother goddess \textbf{Ninmah} (also Ninhursag) shapes humans from \textbf{clay} drawn over the deep. In the \textbf{Atrahasis} text a god is slain, and humanity is made from clay mixed with the \textbf{blood of that god}, so that humans carry a spirit and bear the gods' workload.

\subsubsection{The Babylonian Enuma Elish}

The \textbf{Enuma Elish}, which means \enquote{When on high,} is written on \textbf{7 tablets}. Its purpose is to crown \textbf{Marduk}, the god of Babylon, as king of the gods.

\begin{itemize}
 \item \textbf{Apsu} is the fresh water. \textbf{Tiamat} is the salt water. Their mingling gives birth to the gods.
 \item The young gods grow noisy. Apsu plots to destroy them. \textbf{Ea} (Enki) learns the plot, kills Apsu, and builds his home on Apsu's body.
 \item \textbf{Marduk} is born to Ea. Tiamat breeds an army of monsters and sets \textbf{Kingu} at its head, giving him the \textbf{Tablet of Destinies}.
 \item Marduk is chosen as champion and granted supreme kingship. He slays Tiamat and splits her body to make \textbf{the heavens above and the earth below}. From her body he forms clouds, rivers, and mountains.
 \item Marduk creates \textbf{humankind from the blood of Kingu} to serve the gods. The gods build \textbf{Babylon} and the temple \textbf{Esagila}, and proclaim the \textbf{50 names of Marduk}.
\end{itemize}

The two accounts differ on how the world is made. The Sumerian frame is \textbf{birth from a mother}. The Babylonian frame is \textbf{victory in combat}. The Assyrian version substitutes \textbf{Ashur} for Marduk as the high god.

\subsubsection{The Great Gods}

\begin{longtable}{L{2.2cm} L{2.2cm} L{4.3cm} L{4.3cm}}
\caption{The leading gods across the two languages.}\\
\toprule
\textbf{Sumerian} & \textbf{Akkadian} & \textbf{Domain} & \textbf{Role} \\
\midrule
An & Anu & sky & remote sky-father, source of authority \\
Enlil & Enlil & air, storm, kingship & active king in the Sumerian era \\
Enki & Ea & wisdom, fresh water, craft & orders the world, friend of humans \\
Ki / Ninhursag & Mami & earth, birth & mother goddess, shapes humans \\
Inanna & Ishtar & love and war, Venus & most prominent goddess \\
Nanna & Sin & the moon & father of Utu and Inanna \\
Utu & Shamash & the sun, justice & judge of gods and humans \\
Ishkur & Adad & storm, rain & weather god \\
Marduk & Marduk & kingship, Babylon & victor over Tiamat \\
Ereshkigal & Ereshkigal & the underworld & queen of the dead \\
\bottomrule
\end{longtable}

\subsection{The Anunnaki and the Me}

The gods sit in ranks. The \textbf{Anunnaki} are the great gods, the high assembly, the \textbf{judges} who decree fates. In some texts they become judges of the dead. Below them are the \textbf{Igigi}, the younger laboring gods of heaven. In \textbf{Atrahasis} the Igigi do the digging, then rebel, and this revolt leads to the making of humans to take over the work.

Two instruments govern the whole order.

\begin{itemize}
 \item The divine decrees of civilization (the \textbf{me}) are concrete powers, listed in long catalogs, covering kingship, priesthood, crafts, music, truth, and even falsehood. They make ordered society possible. In \textbf{Inanna and Enki} the goddess takes the me from Enki at Eridu and brings them to her city Uruk.
 \item The \textbf{Tablet of Destinies} is the warrant of supreme rule. Whoever holds it commands the cosmos. Tiamat gives it to Kingu, and Marduk seizes it after his victory.
\end{itemize}

Each person also had a \textbf{personal god}, a guardian who interceded for them before the great gods. When the personal god departed, the person was left open to demons and misfortune.

\subsection{The Sexagesimal God-Numbers}

The tradition runs on a \textbf{base-60} counting system. The great gods each carry a \textbf{sacred number} drawn from that system, forming a descending scale of rank.

\begin{longtable}{L{3.5cm} L{2cm} L{6cm}}
\caption{The sacred numbers of the great gods.}\\
\toprule
\textbf{God} & \textbf{Number} & \textbf{Domain} \\
\midrule
Anu & 60 & sky, highest authority \\
Enlil & 50 & air, kingship \\
Ea (Enki) & 40 & wisdom, fresh water \\
Sin (Nanna) & 30 & moon \\
Shamash (Utu) & 20 & sun, justice \\
Ishtar (Inanna) & 15 & Venus, love and war \\
Adad (Ishkur) & 10 & storm \\
\bottomrule
\end{longtable}

When Marduk rises to kingship he inherits \textbf{Enlil's 50}, which matches the 50 names proclaimed for him. The moon god's \textbf{30} matches the roughly 30 days of the lunar month. The numbered theology is bound to the sexagesimal mathematics that gave us the 60-minute hour and the 360-degree circle.

\subsection{The Underworld and the Descent of Inanna}

Death is grim and without paradise. The dead go down to a dark land below and stay there. The underworld is called \textbf{Kur} (Sumerian) or \textbf{Irkalla} (Akkadian), the \textbf{land of no return}. It is ruled by the queen \textbf{Ereshkigal} and her consort \textbf{Nergal}. The dead eat dust and clay. There is no reward and no sorting of the good from the wicked. All share the same gloom. What an afterlife is like depends less on conduct and more on \textbf{proper burial} and the food and water offerings of the living, the rite called \textbf{kispu}.

The central underworld myth is the \textbf{Descent of Inanna} (the Akkadian \textbf{Descent of Ishtar} is a shorter parallel).

\begin{itemize}
 \item Inanna sets out for the underworld of her sister Ereshkigal and tells her helper \textbf{Ninshubur} to seek rescue if she does not return.
 \item She passes through \textbf{7 gates}. At each gate a garment or ornament is removed, until she stands naked and powerless.
 \item The Anunnaki judges fix her with the eye of death. Inanna dies and her corpse is hung on a hook.
 \item \textbf{Enki} contrives her rescue with the water and food of life.
 \item The underworld demands a \textbf{substitute}. Inanna finds her husband \textbf{Dumuzi} (Akkadian \textbf{Tammuz}) failing to mourn her and gives him over.
 \item A compromise sets up a \textbf{cycle}. Dumuzi spends half the year below and his sister the other half. This is read as the \textbf{dying-and-rising vegetation cycle}, the death and return of growth.
\end{itemize}

\subsection{Human Mortality and Gilgamesh}

The deepest treatment of death is the \textbf{Epic of Gilgamesh}. Gilgamesh, king of Uruk, loses his friend \textbf{Enkidu} and is gripped by the fear of his own death. He seeks immortality.

\begin{itemize}
 \item He reaches \textbf{Utnapishtim}, the one human granted eternal life, who tells the \textbf{flood story}. The gods sent a flood, but Ea warned him to build a boat and save life. This is the same flood found in Atrahasis and the Eridu Genesis, with the survivor named Atrahasis or Ziusudra there.
 \item Utnapishtim challenges Gilgamesh to stay awake for six days and seven nights. He fails at once, showing he cannot conquer even sleep, let alone death.
 \item Gilgamesh wins a plant that restores youth, but a \textbf{snake} steals it while he bathes, taking the renewal for itself.
 \item He returns empty-handed. The lesson is that \textbf{death is the human lot}. Immortality belongs to the gods. What endures is the city and the works a person leaves behind.
\end{itemize}

\subsection{Comparison}

The Mesopotamian model is the precursor of the later traditions of its region. It already holds the watery beginning, the council of gods, the splitting of one body into sky and earth, and the making of humans from clay to serve. It differs from the more hopeful systems in its joyless afterlife, where the going-out has no bright return and what remains is only the work one leaves behind.
